% !TeX program = tectonic
\documentclass[11pt,a4paper]{article}
\usepackage{fontspec}
\usepackage[russian]{babel}
\babelfont{rm}[Language=Default]{Liberation Serif}
\babelfont[Language=Default]{sf}{Carlito}
\babelfont[Language=Default]{tt}{DejaVu Sans Mono}
\usepackage{amsmath,amssymb,amsthm}
\usepackage{geometry}
\geometry{a4paper, top=2.4cm, bottom=2.4cm, left=2.4cm, right=2.4cm}
\usepackage{booktabs,tabularx,enumitem,xcolor,tcolorbox}
\tcbuselibrary{breakable,skins}
\definecolor{linkblue}{HTML}{1F4E79}
\definecolor{statgreen}{HTML}{2F6B3A}
\definecolor{goldacc}{HTML}{8B7E5A}
\usepackage[colorlinks=true,linkcolor=linkblue,urlcolor=linkblue,unicode=true]{hyperref}
\theoremstyle{plain}
\newtheorem{theorem}{Теорема}
\newtheorem{lemma}[theorem]{Лемма}
\newtheorem{proposition}[theorem]{Предложение}
\newtheorem{corollary}[theorem]{Следствие}
\theoremstyle{definition}
\newtheorem{definition}[theorem]{Определение}
\theoremstyle{remark}
\newtheorem{remark}[theorem]{Замечание}
\newtcolorbox{statusbox}[1][]{colback=statgreen!5,colframe=statgreen!75,
  title={\textbf{Сводка доказанного:} #1},fonttitle=\small,boxrule=0.7pt,arc=2pt,breakable}
\newtcolorbox{databox}[1][]{colback=goldacc!6,colframe=goldacc!80,
  title={\textbf{Данные протокола:} #1},fonttitle=\small,boxrule=0.7pt,arc=2pt,breakable}
\newtcolorbox{verifybox}[1][]{colback=linkblue!5,colframe=linkblue!80,
  title={\textbf{Верификация:} #1},fonttitle=\small,boxrule=0.7pt,arc=2pt,breakable}
\widowpenalty=10000
\clubpenalty=10000
\setlength{\parskip}{0.35em}
\newcommand{\CC}{\mathbb{C}}
\newcommand{\RR}{\mathbb{R}}
\newcommand{\ZZ}{\mathbb{Z}}
\newcommand{\QQ}{\mathbb{Q}}
\newcommand{\Ga}{\Gamma}
\newcommand{\im}{\operatorname{Im}}
\newcommand{\re}{\operatorname{Re}}
\newcommand{\disc}{\operatorname{disc}}
\begin{document}
\begin{center}
{\LARGE\bfseries Перепись проводников: две схемы --- один род}\\[0.4em]
{\large Точные гистограммы характеров на четырёх ступенях $N=7/9/15/30$}\\[0.6em]
{\normalsize Исаев Исхак Хамзатович}\\[0.2em]
{\small Программа <<Динамический принцип>> --- лаборатория \texttt{hodge-laboratory}}\\[0.15em]
{\small Отдельная монография теоремы · Монография 20/20}
\end{center}
\vspace{0.4em}
{\color{linkblue}\hrule height 0.8pt}
\vspace{0.8em}

\section{Постановка}

Если монография 1/20 выводила род кривой Ферма топологически
(Риман--Гурвиц), то здесь род получается \emph{арифметически} ---
как итог переписи собственных характеров по проводникам.
Перепись --- арифметический скелет лестницы: она определяет
кратности $h_d$, с которыми проводники $d\mid N$ входят во все
произведения Коула--Селберга программы (сертификаты G и H), и
работает одинаково на всех четырёх ступенях $N=7,9,15,30$.
Протокол V1 исполняет \emph{две независимые схемы} переписи ---
прямую и моёбиусову --- и требует их побитового согласия; эта
монография доказывает, что обе схемы считают одно и то же, и
собирает точные гистограммы всех ступеней.

\begin{definition}[собственный характер и проводник]\label{th:def}
Собственный характер уровня $N$ --- пара $(a,b)$ целых чисел с
$1\le a$, $1\le b$, $a+b\le N-1$; он задаёт голоморфную дифференциальную
форму $\omega_{a,b}=x^ay^b\,dx/y^b$ на кривой Ферма
$X_N\colon x^N+y^N=1$. Проводник характера ---
$d=N/\gcd(N,a,b)$ --- наименьший уровень, через который факторизуется
характер. Перепись --- гистограмма $h_d$ проводников;
$c(k)=(k-1)(k-2)/2$ --- число всех собственных характеров
уровня $k$.
\end{definition}

\begin{theorem}[треугольник]\label{th:tri}
$c(N)=(N-1)(N-2)/2$, и сумма переписи равна роду:
$\sum_{d\mid N}h_d=c(N)=(N-1)(N-2)/2=g$.
\end{theorem}

\begin{proof}
Число пар $1\le a,b$, $a+b\le N-1$ --- точки целочисленного
треугольника: для $a=1,\dots,N-2$ ровно $N-1-a$ значений $b$,
итого $\sum_{a=1}^{N-2}(N-1-a)=1+2+\dots+(N-2)=(N-1)(N-2)/2$.
Каждая пара даёт базисную голоморфную форму (классика для кривой
Ферма, монография 1/20), поэтому $c(N)=g$ --- арифметическая
сторона формулы рода, независимая от Римана--Гурвица. Поскольку
каждый характер имеет ровно один проводник
(лемма~\ref{th:lift} ниже показывает, что проводники делят $N$),
сумма гистограммы и есть $c(N)$.
\end{proof}

\begin{theorem}[лемма подъёма и обращение Мёбиуса]\label{th:lift}
Для каждого $d\mid N$ отображение
$(a',b')\mapsto\big(\tfrac{N}{d}a',\tfrac{N}{d}b'\big)$ --- биекция
между собственными характерами уровня $d$ и собственными
характерами уровня $N$ с проводником, делящим $d$. Следовательно,
\[
  \sum_{d'\mid d}h_{d'}=c(d),
  \qquad\text{и по обращению Мёбиуса}\qquad
  h_d=\sum_{m\mid d}\mu(m)\,c(d/m).
\]
В частности, $h_d$ не зависит от уровня $N$ (лишь $d\mid N$):
$h_d(N)=h_d(d)$ --- \emph{самоподобие лестницы}: каждый проводник
несёт полный меньший стенд.
\end{theorem}

\begin{proof}
\emph{Корректность.} Пусть $(a',b')$ --- характер уровня $d$.
Положим $a=\tfrac{N}{d}a'$, $b=\tfrac{N}{d}b'$: тогда
$1\le a,b$ и $a+b=\tfrac{N}{d}(a'+b')\le\tfrac{N}{d}(d-1)
=N-\tfrac{N}{d}\le N-1$, так что $(a,b)$ --- характер уровня $N$.
Его проводник: $\gcd(N,a,b)=\tfrac{N}{d}\gcd(d,a',b')$, поэтому
$d(a,b)=N/\gcd(N,a,b)=d/\gcd(d,a',b')=d(a',b')\mid d$.

\emph{Сюръективность с разложением.} Если характер $(a,b)$
уровня $N$ имеет проводник $d'\mid d$, то
$\gcd(N,a,b)=\tfrac{N}{d'}$ кратно $\tfrac{N}{d}$, откуда
$\tfrac{N}{d}\mid a,b$, и пара $a'=\tfrac{d}{N}a$,
$b'=\tfrac{d}{N}b$ целая; границы сохраняются обратным
рассуждением корректности, и
корректности, и $(a',b')$ --- характер уровня $d$ с проводником
$d'$, восстанавливающий $(a,b)$. Инъективность очевидна. Значит
$\sum_{d'\mid d}h_{d'}(N)=c(d)$ --- число всех характеров
уровня $d$.

\emph{Обращение.} Решётка делителей $d$ --- конечный
дистрибутивная решётка, и функция Мёбиуса на ней --- классическая
$\mu$; обращение Мёбиуса на решётке делителей даёт вторую
формулу. Правая часть зависит только от $d$, поэтому $h_d$
одинаков на всех уровнях $N$, кратных $d$: в уровень $d$ формула
даёт $h_d(d)$. Самоподобие доказано.
\end{proof}

\begin{theorem}[переписи четырёх ступеней]\label{th:four}
Точные гистограммы проводников:
\[
\begin{array}{lll}
N=7: & \{7\colon 15\}, & 15=15;\\[0.15em]
N=9: & \{3\colon 1,\ 9\colon 27\}, & 28=1+27;\\[0.15em]
N=15: & \{3\colon 1,\ 5\colon 6,\ 15\colon 84\}, & 91=1+6+84;\\[0.15em]
N=30: & \{3\colon 1,\ 5\colon 6,\ 6\colon 9,\ 10\colon 30,\ 15\colon 84,\ 30\colon 276\}, & 406=1+6+9+30+84+276.
\end{array}
\]
Каждая строка проверяется обеими схемами, включая явные значения
$h_6=c(6)-c(3)-c(2)+c(1)=10-1-0+0=9$ и
$h_{10}=c(10)-c(5)-c(2)+c(1)=36-6=30$,
$h_{30}=c(30)-c(15)-c(10)-c(6)+c(5)+c(3)+c(2)-c(1)
=406-91-36-10+6+1=276$.
\end{theorem}

\begin{proof}
\emph{Прямая схема.} $N=7$ прост: $\gcd(7,a,b)=1$ для всех
характеров треугольника --- единственный проводник $7$,
$h_7=c(7)=15$. $N=9$: проводник $3$ означает
$\gcd(9,a,b)=3$, т.е. $a,b\in\{3,6\}$, и $a+b\le8$ оставляет
$(3,3)$: $h_3=1$; остаток $h_9=28-1=27$. $N=15$: проводник $3$
--- $\gcd(15,a,b)=5$, т.е. $a,b\in\{5,10\}$, $a+b\le14$ ---
$(5,5)$: $h_3=1$; проводник $5$ --- $\gcd=3$, т.е. $a,b\in\{3,6,9,12\}$
с $a+b\le14$ --- шесть пар $(3,3),(3,6),(6,3),(3,9),(9,3),(6,6)$:
$h_5=6$; остаток $h_{15}=91-1-6=84$. $N=30$ --- проводник $6$: $\gcd(30,a,b)=5$ --- $a,b$ кратны $5$,
$a,b\in\{5,10,15,20,25\}$ с $a+b\le29$ и $\gcd(30,a,b)=5$
(не $10$, не $15$) --- девять пар; моёбиусова схема даёт то же
$h_6=9$; аналогично $h_{10}=30$, $h_{15}=84$ (самоподобие:
$h_{15}(30)=h_{15}(15)$), $h_{30}=276$.

\emph{Моёбиусова схема.} Формула теоремы~\ref{th:lift} с
явными значениями $c$: $c(2)=0$, $c(3)=1$, $c(5)=6$, $c(6)=10$,
$c(9)=28$, $c(10)=36$, $c(15)=91$, $c(30)=406$; вычисления
приведены в формулировке --- все совпадают с прямой схемой
попарно по каждому проводнику.
\end{proof}

\begin{theorem}[две схемы --- одно целое]\label{th:two}
Прямая схема (перебор пар с вычислением $d=N/\gcd(N,a,b)$) и
моёбиусова схема (включения-исключения по делителям) дают
\emph{побитово одинаковые} гистограммы на всех уровнях
$N=7,9,15,30$ --- целочисленное тождество, детерминированная
проверка без порогов.
\end{theorem}

\begin{proof}
По теореме~\ref{th:lift} обе схемы вычисляют одну и ту же
функцию $h_d$: прямая --- по определению (каждая пара попадает
ровно в один проводник), моёбиусова --- как единственное решение
уравнения свёртки $\sum_{d'\mid d}h_{d'}=c(d)$ на решётке
делителей. Формальное совпадение не оставляет свободы
расхождениям; машинное исполнение (протокол V1) --- отпечаток
против ошибок реализации: два независимых кодовых пути обязаны
произвести одинаковые словари $\{d\colon h_d\}$, что и
фиксируется побитовым сравнением.
\end{proof}

\begin{remark}
Гистограммы экспоненциально сжаты квадратичным характером
уровня: у простого $7$ --- один проводник; у $9=3^2$ --- два
(появляется <<отшельник>> $(3,3)$ --- след уровня $3$); у
$15=3\cdot5$ --- три; у $30=2\cdot3\cdot5$ --- шесть: все делители уровня,
кроме $1$ и $2$ (треугольник уровня $d$ пуст лишь при $d\le2$). Самоподобие
$h_d(N)=h_d(d)$ превращает лестницу в фрактальное семейство
стендов: большая ступень содержит все меньшие как проводниковые
слои, и каждое слагаемое суммы $406=1+6+9+30+84+276$ само есть
род меньшего стенда (за вычетом его собственных меньших
проводников). Именно кратности $h_d$ входят в показатели
произведения Коула--Селберга сертификата H --- перепись замыкает
$\Ga$-архитектуру программы.
\end{remark}

\begin{databox}[Перепись --- сводка четырёх ступеней]
$N=7$: род $15$, $\{7\colon15\}$, один проводник (простой
уровень). $N=9$: род $28$, $\{3\colon1,\,9\colon27\}$,
диагональный характер $(3,3)$. $N=15$: род $91$,
$\{3\colon1,\,5\colon6,\,15\colon84\}$. $N=30$: род $406$,
$\{3\colon1,\,5\colon6,\,6\colon9,\,10\colon30,\,15\colon84,\,30\colon276\}$.
Все восемь строк подтверждены двумя схемами (V1); суммы равны
$(N-1)(N-2)/2$ точно; самоподобие $h_d(N)=h_d(d)$ (например,
$h_{15}(30)=84=h_{15}(15)$). Эталон: \texttt{baseline\_v1\_v9.json},
блок \texttt{V1\_census}.
\end{databox}

\begin{statusbox}[итог]
Арифметика рода поставлена на собственный фундамент: треугольник
характеров, лемма подъёма с биекцией проводников, обращение
Мёбиуса, самоподобие $h_d(N)=h_d(d)$, точные гистограммы всех
четырёх ступеней и побитовое согласие двух схем. Род ---
считается, а не постулируется; кратности $h_d$ --- вход
$\Ga$-нормировки сертификатов G/H.
\end{statusbox}

\begin{verifybox}[как воспроизвести]
\texttt{python3 laboratory.py} $\to$ протокол V1 (перепись,
две схемы, уровни $7/9/15/30$); \texttt{--check-baseline}
сверяет гистограммы и роды с эталоном; Lean: блок
\texttt{genus\_fermat}; batch-режим: пункт \texttt{census}
в сценарии.
\end{verifybox}

\vfill
{\color{linkblue}\hrule height 0.6pt}
\vspace{0.5em}
{\small\color{gray}
Лицензия: индивидуальная исключительная лицензия Исаева Исхака Хамзатовича.
Все права на вычисления, формулы и тексты принадлежат автору.
Верификация: \texttt{python3 laboratory.py} (протокол V1, \texttt{--check-baseline}).
}
\end{document}
